\section{Methods}
% \danny{Start rough draft of supplementary}
\vspace{-1ex}

\begin{figure*}[t]
\centering
    \includegraphics[width=\textwidth]{figures/schematic_big.pdf}
    \vspace{-10pt}
    \caption{\textbf{CoarsenConf architecture.} \textbf{(I)} The encoder $q_\phi(z| X, \gR)$ takes ground truth conformer $X$, RDKit approximate conformer $\gR$, and coarse-grained (CG) conformer $\gC$ as inputs (derived from $X$ and predefined CG strategy), and outputs a variable-length equivariant CG representation via equivariant message passing and point convolutions. 
    \textbf{(II)} Equivariant MLPs are applied to learn the mean and log variance of both the posterior and prior distributions. 
    \textbf{(III)} The posterior (training) or prior (inference) is sampled and fed into the Channel Selection module, where an attention layer is used to learn the optimal pathway from CG to FG structure. 
    \textbf{(IV)} Given the FG latent vector and the RDKit approximation, the decoder $p_\theta(X |\gR,z)$ learns to recover the low-energy FG structure through autoregressive equivariant message passing. 
    The entire model can be trained end-to-end by optimizing the KL divergence of latent distributions and reconstruction error of generated conformers.
    }
    \label{fig:schematic}
    \vspace{-5pt}
\end{figure*}
\subsection{Algorithmic Details}
\label{sec:methods-algorithmic-details}

\vspace{-1ex}
\paragraph{Coarse-graining Procedure.}
We first define a rotatable bond as any single bond between two non-terminal atoms, excluding amides and conjugated double bonds, where the torsion angle is the angle of rotation around the central bond. Formally, the torsion angle $\tau_{abcd}$ is defined about bond $(b, c) \in \gE$ where (a, b) are a choice of reference neighbors $\text{s.t } a \in \mathcal{N}(b)\setminus c \text{ and } d \in \mathcal{N}(c)\setminus b$. 

We coarsen molecules into a single bead for each connected component resulting from severing all rotatable bonds.
This choice in CG procedure implicitly forces the model to learn over torsion angles as well as atomic coordinates and inter-atomic distances.
We found that using a more physically constrained definition of torsional angles, as defined by \citet{GeoMol}, in the CG procedure led to a significant increase in performance compared to that used in \citet{torsional-diffusion}. This is because the latter allows rotations around double and triple bonds, while the former does not. 
An example of the coarse-graining procedure is in~\fref{fig:novelty}(I). For formal definitions, see Appendix \sref{section:appendix-CG}.

\vspace{-1ex}
\paragraph{Learning Framework.} \model is a conditional generative model that learns $p(X | \gR)$ where $X$ is the low-energy 3D conformation, and $\gR$ is the RDKit best guess conformation. Specifically, we optimize $p(X | \gR)$ by maximizing its variational lower bound with an approximate posterior distribution $q_\phi(z |X, \gR)$ and learned prior $p_\psi(z | \gR)$:
\begin{equation}
\begin{aligned}
    \log p(X|\gR) \geq & \underbrace{\E_{q_\phi(z | X, \gR)} \log p_\theta(X | \gR, z)}_{\gL_{\text{reconstruction}}} + \underbrace{\E_{q_\phi(z | X, \gR)} \log \frac{ p_\psi(z|\gR) }{q_\phi(z| X, \gR)}}_{\gL_{\text{latent regularization}}}, %:= \gL_{\text{ELBO}},
\label{eq:elbo}
\end{aligned}
\end{equation}
where $q_\phi(z|X,\gR)$ is the hierarchical equivariant encoder model, $p_\theta(X|\gR,z)$ is the equivariant decoder model to recover $X$ from $\gR$ and $z$, and $p_\psi(z|\gR)$ is the learned prior distribution. The reconstruction loss, $\gL_{\text{recon.}}$, is implemented as $\text{MSE}(\gA(X_{true}, X_{model}), X_{model})$, where $\gA$ is the Kabsch alignment function that provides an optimal rotation matrix and translation vector to minimize the mean squared error (MSE) \citep{Kabsch}. The second term, $\gL_{\text{reg.}}$, can be viewed as a regularization over the latent space and is implemented as $\beta D_{KL}(q_\phi(z | X, \gR)\parallel p_\psi(z|\gR))$ \citep{beta-vae}. In practice, we add an auxiliary loss function that measures the distance error between the coordinates of all 2-hop neighbors in the molecular graph. More details on the loss function are in Appendix~\sref{section:appendix-loss-func}.

\vspace{-1ex}
\paragraph{Learning Optimal Distortions.} As the goal of MCG is to model the conditional distribution $p(X | \gR)$, we can simplify the learning objective by setting $X = \gR + \Delta X$ and learn the optimal distortion $\Delta X$ from the RDKit approximation. 
This greatly simplifies the learning objective, as we have ensured $\gL_{\text{recon.}}$ is no worse than that of the RDKit approximation.
This procedure is analogous to using multiple sequence alignments in protein structure prediction tasks to initialize the model in an approximate low-energy state \citep{alphafold}. We stress that the RDKit approximate conformers are trivial to obtain compared to the cost of generative model training and inference. 

\vspace{-1ex}
\subsection{Encoder Architecture}
\label{section:hierarchy}
\paragraph{Hierarchical Modules.} We describe the encoder, shown in~\fref{fig:schematic}(I). The model operates over SE(3)-invariant atom features $h \in R^{n\times D}$, and SE(3)-equivariant atomistic coordinates $x \in R^{n\times3}$. A single encoder layer is composed of three modules: fine-grained, pooling, and coarse-grained. Full equations for each module can be found in Appendix \sref{section:eqn-fine}, \sref{section:eqn-pool}, \sref{section:eqn-coarse}, respectively.

The fine-grained module is a graph-matching message-passing architecture. It differs from~\citet{equibind} by not having internal closed-form distance regularization and exclusively using unidirectional attention. Its purpose is to effectively match the approximate conformer and ground truth by updating attention from the former to the latter.

The pooling module takes in the updated representations ($h$ and $x$) of both the ground truth molecule and the RDKit reference from the FG module.
The pooling module is similar to the FG module, except it no longer uses attention and operates over a pooling graph.
Given a molecule with $n$ atoms and $N$ CG beads, the pooling graph consists of $n + N$ nodes. There is a single directional edge from all atoms to their respective beads (see~\fref{fig:novelty}(I)). This allows message passing to propagate information through the predefined coarsening strategy.

The coarse-grained module uses the updated CG representations ($H\in R^{N\times D}$ and $X \in R^{N\times 3}$) from the pooling module to learn equivariant CG features ($Z$ and $\Tilde{Z} \in R^{N\times F \times 3}$) for the ground truth molecule and the RDKit reference. $F$ is fixed as a hyperparameter for latent space size. $N$ is allowed to be variable-length to handle molecules resulting from any coarsening procedure. 
The CG features are learned using a graph-matching point convolution~\citep{tensor-field-network} with similar unidirectional attention as the FG module. Prior to the main message-passing operations, the input features undergo equivariant mixing~\citep{3dlinker} to further distill geometric information into the learned CG representation.

The modules in each encoder layer communicate with the respective module of the previous layer. This hierarchical message-passing scheme results in an informative and geometrically grounded final CG latent representation. We note that the pooling module of layer $\ell$ uses the updated invariant features $H$ from the CG module of layer $\ell - 1$, as shown by the dashed lines in~\fref{fig:schematic}(I).

\begin{figure}[t]
\centering
    \includegraphics[width=\textwidth]{figures/attention.png}
    \vspace{-10pt}
    \caption{\textbf{Variable length coarse-to-fine backmapping via Aggregated Attention.} The latent vector of a single CG bead $Z_{B}$ $\in R^{F \times 3}$ is used as the key and value of a single head attention operation with an embedding dimension of three to match the x,y,z coordinates. The query vector is the subset of the RDKit conformer corresponding to bead $B$ $\in R^{n_{B} \times 3}$, where $n_B$ is variable-length as we know a priori how many FG atoms correspond to a certain CG bead. 
    Leveraging attention, we efficiently learn the optimal blending of latent features for FG reconstruction. We call this approach Aggregated Attention, as it aggregated 3D segments of FG information to form our latent query.
     }
    \label{fig:attention}
    \vspace{-5pt}
\end{figure}

\vspace{-1ex}
\paragraph{Equivariant Latent Space.} As $Z$ holds a mixture of equivariant spatial information, we maintain equivariance through the reparametrization trick of the VAE (\fref{fig:schematic}(II)). Specifically, we define the posterior and prior means ($\bm{\mu}_\phi$, $\bm{\mu}_\psi$) and standard deviations ($\bm{\sigma}_\phi$, $\bm{\sigma}_\psi$), as follows:
\begin{equation} 
\begin{aligned}
\bm{\mu}_\phi &= \text{VN-MLP}(Z, \Tilde{Z}), \\
\bm{\mu}_\psi &= \text{VN-MLP}(\Tilde{Z}), \\
\end{aligned}
\quad
\begin{aligned}
\bm{\text{log}(\sigma^2}_\phi) &= \text{MLP}(Z, \Tilde{Z}), \\
\bm{\text{log}(\sigma^2}_\psi) &= \text{MLP}(\Tilde{Z}).
\end{aligned}
\end{equation}
We use an invariant MLP to learn the variance and apply it to the x, y, and z directions to enforce equivariance. We found that setting $\sigma_\psi = 1$ results in smoother optimization rather than requiring an additional latent regularization to the standard normal as seen in \citet{CVGAE}.
We note the conditional posterior is parameterized with both the ground truth and RDKit approximation, whereas the learned conditional prior only uses the approximation.
\vspace{-1ex}
\subsection{Decoder Architecture}
\vspace{-1ex}

We sample from the learned posterior (training) and learned prior (inference) to get $Z = \mu + \eps \sigma$, where $\eps$ is noise sampled from a standard Gaussian distribution as the input to the decoder. 
We note the role of the decoder is two-fold. The first is to convert the latent coarsened representation back into FG space through a process we call channel selection. The second is to refine the fine-grain representation autoregressively to generate the final low-energy coordinates.

\vspace{-1ex}
\paragraph{Channel Selection.} To explicitly handle all choices of coarse-graining techniques, our model performs variable-length backmapping.
% To improve conformer generation using coarse-graining, we need to ensure that our model can handle variable-length backmapping. 
This aspect is crucial because every molecule can be coarsened into a different number of beads, and there is no explicit limit to the number of atoms a single bead can represent. Unlike CGVAE~\citep{CGVAE}, which requires training a separate model for each choice in granularity $N$, \model is capable of reconstructing FG coordinates from any $N$ (illustrated in \fref{fig:schematic}(III)).

CGVAE defines the process of channel selection as selecting the top $k$ latent channels, where $k$ is the number of atoms in a CG bead of interest. Instead of discarding all learned information in the remaining $F-k$ channels in the latent representation, we use a novel aggregated attention mechanism. This mechanism learns the optimal mixing of channels to reconstruct the FG coordinates and is illustrated in~\fref{fig:attention}.
The attention operation allows us to actively query our latent representation for the number of atoms we need, and draw upon similarities to the learned RDKit approximation that has been distilled into the latent space through the encoding process.
Channel selection translates the CG latent tensor $Z \in R^{N\times F \times 3}$ into FG coordinates $X_{CC} \in R^{n\times 3}$.

\vspace{-1ex}
\paragraph{Autoregressive Refinement.} Once channel selection is complete, we have effectively translated the variable-length CG representation back into the desired FG form. From here, $X_{CC}$ is grouped into its corresponding CG beads, but left in FG coordinates to do a bead-wise autoregressive generation of final low-energy coordinates (\fref{fig:schematic}(IV)). 
As there is no intrinsic ordering of subgraphs, we use a breadth-first search that prioritizes larger subgraphs with large out-degrees. In other words, we generate a linear order that focuses on the largest, most connected subgraphs and works outward. We believe that by focusing on the most central component first, which occupies the most 3D volume, we can reduce the propagation of error that is typically observed in autoregressive approaches.
We stress that by coarse-graining by torsion angle connectivity, our model learns the optimal torsion angles in an unsupervised manner as the conditional input to the decoder is not aligned. \model ensures each next generated subgraph is rotated properly to achieve a low coordinate and distance error.

\vspace{-1ex}
\paragraph{Learning the Optimal Distortion.}The decoder architecture is similar to the EGNN-based FG layer in the encoder. However, it differs in two important ways. First, we mix the conditional coordinates with the invariant atom features using a similar procedure as in the CG layer instead of typical graph matching. Second, we learn to predict the difference between the RDKit reference and ground truth conformations. This provides an upper error bound and also enables us to leverage easy-to-obtain approximations more effectively.
More formally, a single decoder layer is defined as follows:
\begin{subequations}
\begin{align}
\bm{\mu}^{(t)} &= \frac{1}{|\gV_{prev}|} \sum_{k\in \gV_{prev}}x_k, \\
\bm{\Tilde{h}}_{i} &= \phi^m(\bm{h}_i^{(t)},\bm{x}_i^{(t)},\bm{\mu}^{(t)},\norm{\bm{x}_i^{(t)}-\bm{\mu}^{(t)}}^2),\forall i \in \gV_{cur}, \\
\bm{m}_{j\to i} &= \phi^e(\bm{\Tilde{h}}_i^{(t)},\bm{\Tilde{h}}_j^{(t)},\norm{\bm{x}_i^{(t)}-\bm{x}_j^{(t)}}^2,\norm{\bm{x}_i^{(t)}-\bm{x}_{ref, j}^{(t)}}^2, \norm{\bm{x}_i^{(t)}-\bm{x}_{ref, i}^{(t)}}^2),\forall (i,j)\in \gE_{cur}, \\
\bm{m}_i &= \frac{1}{|\gN(i)|} \sum_{j\in\gN(i)}\bm{m}_{j\to i},\forall i\in\gV_{cur}, \\
\bm{u}_{j'\to i} &= a_{j'\to i} \bm{W} \bm{h}_{j'}^{(t)},\forall i\in\gV_{cur},j'\in\gV_{prev}, \\
\bm{u}_{i} &=  \sum_{j'\in\gV_{prev}}\bm{u}_{j'\to i},\forall i \in\gV_{cur},\\
\bm{x}_i^{(t+1)} &= \bm{x}_{ref,i}^{(t)} + \sum_{j\in\gN(i)} (\bm{x}_i^{(t)}-\bm{x}_j^{(t)})\phi^x(\bm{m}_{j\to i}), \\
\bm{h}_i^{(t+1)} &= (1 - \beta) \cdot \bm{h}_i^{(t)} + \beta \cdot \phi^h(\bm{\Tilde{h}}_i^{(t)},\bm{m}_i,\bm{u}_{i}, \bm{f}_i),\forall i\in\gV_{cur},
\end{align}
\label{eqn:decoder}
\end{subequations}
where ($\gV_{cur}$, $\gE_{cur}$) and ($\gV_{prev}$, $\gE_{prev}$) refer to the subgraph currently being generated and the set of all previously generated subgraphs, \ie, the current state of the molecule. $\phi^m, \phi^e, \phi^x, \text{and } \phi^h$ refer to separate shallow MLPs for the feature mixing, edge message calculation, coordinate update, and invariant feature update, respectively. \eref{eqn:decoder}(a-b) creates a mixed feature for each atom comprised of the current FG invariant feature and 3D position vectors ($h$ and $x$), and the previous centroid $\mu$ and respective centroid distances. \eref{eqn:decoder}(c-d) defines the message passing operation that uses the aforementioned mixed features $\Tilde{h}$ and a series of important distances between the model-based conformer and RDKit reference (see~ \fref{fig:novelty}(II) for a helpful visualization). \eref{eqn:decoder}(e-f) apply the same unidirectional attention updates seen in the encoder architecture. \eref{eqn:decoder}(g-h) update the position and feature vector for each atom using the above messages and attention coefficients, with $f$ representing the original invariant node features $h^{\ell=0}$ and $\beta$ a weighted update parameter. We emphasize that~\eref{eqn:decoder}(g) formulates the overall objective as learning the optimal distortion of the RDKit reference to achieve the low-energy position \ie\ $x^* = x_{ref} + \Delta x$.
The CG autoregressive strategy allows \model to handle extremely large molecules efficiently, as the max number of time steps is equal to the max number of CG beads. \model is trained using teacher forcing~\citep{teacher-forcing}, which enables an explicit mixing of low-energy coordinates with the current FG positions from channel selection \eref{eqn:decoder}(a-b).

\vspace{-2ex}